\documentclass[../main]{subfiles}
\begin{document}
\ifSubfilesClassLoaded{\setcounter{page}{4}}{}
\section{La proprietà comune (pubblica)}
\label{sec:05_proprieta_comune}

L’ostacolo principale che si frappone al progresso verso il comunismo risiede nella proprietà privata. Le riflessioni più mature e articolate di Marx sul comunismo si trovano nei “Manoscritti economico-filosofici del 1844”. Il capitolo dedicato al comunismo è incluso nel secondo manoscritto, immediatamente dopo quello incentrato sulla proprietà privata. Il testo si apre con un’analisi della proprietà privata, del lavoro in quanto sua essenza, e con una disamina del modo in cui la proprietà privata e il lavoro vengono interpretati dai diversi esponenti del socialismo e del comunismo, intesi come correnti ideologiche. I comunisti, già prima di Marx, si oppongono alla proprietà privata dei mezzi di produzione, poiché ritengono che essa assoggetti gli uomini al dominio di altri e generi tutti i mali insiti nella società capitalista.

La proprietà privata e la proprietà personale non costituiscono la stessa categoria. Gli oggetti destinati a un uso immediato e di consumo, appartenenti ai rispettivi possessori, rientrano nella sfera della proprietà privata, così come rientrano in essa fabbriche, stabilimenti, negozi, imprese e vaste estensioni di terreno. Questa distinzione, tuttavia, tende a confondere le acque e ad oscurare la vera natura della questione. Pertanto, quando gli oppositori del comunismo si sentono parlare dell’abolizione della proprietà privata, i più sprovveduti tra loro immaginano una collettivizzazione generalizzata di tutti i beni e si allarmano per la sorte dei propri averi personali. Su questo punto, possiamo rassicurarli. Gli oggetti destinati al consumo personale non rappresentano un interesse per i comunisti nel senso di privare gli individui di tali beni, ma piuttosto nel senso di garantire a ciascuno la possibilità di disporne in quantità sufficiente per condurre una vita piena e appagante in ogni suo aspetto. Dal punto di vista economico, questi oggetti non rientrano nella categoria della proprietà privata, bensì in quella della proprietà personale.

L’abolizione della proprietà privata non implica l’eliminazione del possesso personale – la proprietà privata è, in fin dei conti, un rapporto tra individui e gli oggetti, non una caratteristica intrinseca degli oggetti stessi – poiché il consumo, per sua natura, deve rimanere un atto individuale. Piuttosto, l’abolizione della proprietà privata mira a garantire a tutti i membri della società l’accesso ai beni necessari per il consumo personale e, ove possibile, a trasformare il consumo in forme collettive, più efficienti e razionali. Ad esempio, la creazione di cooperative di consumatori consente di ottimizzare l’utilizzo dei beni, riducendo la necessità di produrne e immagazzinarne in quantità eccessive, e quindi diminuendo i costi in termini di tempo sociale complessivo. In questo modo, anche la semplice istituzione di una cooperativa per la condivisione di beni non utilizzati frequentemente può contribuire a ridurre l’ingombro domestico e a massimizzare l’uso di ciascun oggetto, prolungandone la vita utile sia in termini fisici che di obsolescenza percepita.In tal modo, l’oggetto circola tra i membri della cooperativa, offrendo un vantaggio a ciascuno di essi. Numerosi esempi di cooperative di consumo sono già stati descritti in letteratura. È agevole estendere questo principio a una vasta gamma di beni, che cessano di appartenere a un singolo individuo e vengono utilizzati in modo appropriato da tutti, al fine di soddisfare le loro necessità. In questo contesto, il comunismo si contrappone alla proprietà privata non nella maniera temuta dai suoi detrattori. Le forme di consumo collettivo non mirano a privare alcuno dei propri averi, bensì a garantire a tutti l’accesso ai beni di consumo.

In relazione alle forme di vita comunitarie, quali la condivisione di lavatrici, la preparazione e la distribuzione collettiva dei pasti, l’uso condiviso di strumenti per le pulizie, l’efficienza nell’impiego delle risorse e il risparmio di tempo – sia immediato che mediato, come ad esempio il tempo risparmiato nella produzione di tali risorse – possono rivelarsi significativi, mantenendo invariato il livello di consumo. In ogni circostanza, tale risparmio è proporzionale alle potenzialità offerte dalla tecnologia. Anche oggi, per esempio, preparare l’alimentazione per un bambino in casa e preparare la stessa per duecento bambini in una struttura come un asilo nido o una cucina per neonati, dotate di attrezzature moderne, richiede un lasso di tempo analogo.

Lo stesso principio si applica alla preparazione dei pasti per gli adulti. La ristorazione collettiva permette l'impiego di tecnologie più avanzate, inattuabili nell'ambito della cucina domestica individuale, e consente un grado di automazione decisamente superiore in tutte le fasi del processo. Inoltre, è possibile esercitare un controllo più rigoroso sulla qualità, con benefici che si estendono a tutti coloro che sono coinvolti. Anche in un sistema capitalistico, si riscontra una tendenza analoga. Il capitalismo stesso, a suo modo, sviluppa e standardizza modelli di ristorazione (si pensi, ad esempio, a McDonald's, che offre lo stesso prodotto in ogni parte del mondo), così come le lavanderie, e via dicendo. Tuttavia, in un sistema capitalistico, queste strutture non sono comunque utilizzate in forma pienamente collettiva. Le modalità di consumo collettivo si scontrano costantemente con resistenze all'interno del sistema di produzione, e, come ben sappiamo, è la produzione a determinare i modelli di consumo.

Le innumerevoli difficoltà che affliggono la vita di tutti i giorni potrebbero cedere il passo a soluzioni innovative, fondate su tecnologie avanzate. I veicoli privati, fonte di intralcio e di congestione, lascerebbero il posto a un sistema di trasporto pubblico efficiente e capillare, e così via, in ogni ambito dell'esistenza umana, a beneficio di tutti. Nelle metropoli moderne, ad esempio, dove il trasporto pubblico è ben organizzato, chi usufruisce di questo servizio, soprattutto nelle ore di punta, raggiunge la propria destinazione con maggiore rapidità. A Mosca, i bambini che si recano a scuola in metropolitana giungono puntuali, mentre coloro che vengono accompagnati in auto tendono a subire ritardi, poiché l'uso dello spazio urbano limitato per scopi privati si rivela meno efficace rispetto all'utilizzo di un sistema pubblico ben strutturato.Questo spiega anche perché, in caso di nevicate, le autorità cittadine raccomandano di limitare l'uso dei veicoli privati e di privilegiare il trasporto pubblico, con l'obiettivo di ridurre al minimo il rischio di incidenti stradali.

La produzione di beni incentiva forme di consumo individuali. Pur essendo spesso irrazionali e inefficienti, queste costituiscono una condizione imprescindibile per sostenere una domanda costante di beni sul mercato e per assicurare la circolazione del capitale. In tal senso, la proprietà privata si intreccia con la sfera personale, esercitando su di essa una notevole influenza.

La proprietà privata, o, per essere più precisi, la proprietà privata dei mezzi di produzione, consiste primariamente nella possibilità di ottenere un profitto aggiuntivo (plusvalore), sotto forma di guadagno o di rendita, ovvero di sfruttare la forza lavoro altrui.

I comunisti sostengono la socializzazione della proprietà privata, intendendo con ciò non i beni personali, bensì i mezzi di produzione che, attualmente, anziché essere impiegati per il beneficio dell’intera comunità, servono a perpetuare lo sfruttamento tra gli esseri umani. La proprietà privata non risiede negli oggetti materiali, quanto piuttosto nei rapporti sociali che si generano attraverso di essi; la sua vera natura è lo sfruttamento.

I fautori del comunismo precedenti a Marx concepivano l’abolizione della proprietà privata come un’estensione di quest’ultima all’intera società, trasformando ogni membro in un lavoratore e uniformando le condizioni di tutti, sia in termini di attività lavorativa sia in relazione ai consumi. Per molti sostenitori del comunismo pre-marxista, la proprietà privata era semplicemente intesa come ricchezza materiale, di cui tutti avrebbero potuto essere compartecipi. Ad esempio, Proudhon la considerava unicamente come capitale, destinato all’estinzione. Più vicini alle idee di Marx furono quegli studiosi che individuarono nell’organizzazione del lavoro, in particolare in un lavoro «parcellizzato, frammentato e quindi alienante», la radice della proprietà privata e della sua nefasta influenza sull’individuo, sebbene circoscrissero tale analisi a specifiche categorie di attività.«Come i fisiocrati, Fourier rivaluta l’attività agricola, considerandola la forma di lavoro più elevata; Saint-Simon, al contrario, ritiene che il fulcro della questione risieda nel lavoro industriale e, di conseguenza, auspicava il predominio incontrastato degli industriali e il miglioramento delle condizioni dei lavoratori. Infine, il comunismo rappresenta l’affermazione concreta dell’abolizione della proprietà privata, manifestandosi, nelle sue prime fasi, come una proprietà privata universale». In questo contesto, l’aggettivo «positivo» implica «la salvaguardia di tutte le ricchezze culturali e delle capacità produttive sviluppate dall’uomo nella società mercantile»

Marx sottopone a una critica serrata l'idea di intendere l'abolizione della proprietà privata come sua estensione all'intera società. Egli argomenta che un tale comunismo non supera la proprietà privata, ma ne rappresenta piuttosto la sua forma teorica più compiuta in quanto rapporto sociale. Sottolinea come, in questo comunismo, inteso sia come posizione ideologica e teorica sia come movimento sociale, il dominio della proprietà – concepita come insieme di oggetti – si mantenga ancora inalterato. Le relazioni sociali continuano a essere concepite unicamente in termini di proprietà privata. Un simile comunismo, infatti, "tende a distruggere tutto ciò che non può essere oggetto di appropriazione individuale secondo i principi della proprietà privata, cercando di astrarre in modo forzato da talenti e capacità individuali".A suo avviso, il possesso materiale e immediato costituisce l’unico fine della vita e dell’esistenza; la figura del lavoratore non viene soppressa, bensì estesa a tutti gli esseri umani; il concetto di proprietà privata permane quale rapporto che lega l’intera società al mondo degli oggetti; infine, questo movimento, che si prefigge di opporre alla proprietà privata l’idea di una proprietà privata universale, assume una forma del tutto animalesca quando si contrappone al matrimonio (che, in fondo, rappresenta una specifica forma di proprietà privata esclusiva) con la comunanza delle donne, trasformando queste ultime in una proprietà pubblica e universale.

Secondo Marx, l’idea di una comunanza tra le donne rappresenta una forma di oggettivazione dell’individuo, che riflette, al contempo, il rapporto dell’uomo con il mondo degli oggetti, un mondo creato dall’uomo per soddisfare le esigenze umane. «Così come la donna, in questa prospettiva, passa dal matrimonio alla prostituzione universale, così anche l’intero mondo della ricchezza – ovvero l’essenza materiale dell’uomo – si evolve da un rapporto esclusivo con il proprietario privato a una relazione diffusa con l’intera società». Marx definisce questa concezione di comunismo come «elementare» o «rudimentale». In questo contesto, per ricchezza si intende l’insieme di tutti gli oggetti creati dagli esseri umani, oggetti che incarnano sia le capacità produttive della società umana sia i bisogni che essa ha sviluppato.

La teoria del comunismo rozzo, negando l’individualità, non fa altro che esasperare, nella sua rappresentazione ideologica, la reale tendenza della società verso la privatizzazione della ricchezza. In un simile contesto, dove la trama delle relazioni umane si restringe a rapporti di proprietà, l’individuo è ridotto a mero strumento economico al servizio della proprietà privata. In sostanza, la personalità di tali individui, siano essi proletari o capitalisti, diviene del tutto secondaria rispetto al funzionamento del sistema proprietario. Il proletario interessa al capitalista unicamente come fornitore di una determinata quantità di forza lavoro, e non in quanto persona; analogamente, il capitalista interessa al proletario esclusivamente come acquirente di tale forza lavoro. Lo stesso principio si applica a qualsiasi transazione di compravendita. Ogni altra caratteristica individuale degli agenti economici perde importanza nel contesto del sistema di proprietà. Se l’esistenza di un individuo è interamente assorbita da tali rapporti, la sua personalità finisce per identificarsi esclusivamente con essi.

Il comunismo più elementare postula un’abolizione della proprietà privata che, in realtà, non la intacca realmente. Ciò si manifesta nel suo astratto rigetto delle conquiste culturali, giudicate superflue solo perché ritenute non indispensabili per il proletario in quanto tale, e – come Marx espresse – nella sua tendenza a un’innaturale semplificazione, a una condizione umana ridotta all’essenziale, povera, rozza e priva di aspirazioni, incapace di trascendere la proprietà privata perché non ancora matura per farlo.

Lo slogan “Espropriare e dividere”, che caratterizza i moderni esponenti del cosiddetto comunismo, è ancor più primitivo. Essi non sono nemmeno giunti alla fase dell’espropriazione nella forma in cui era concepita dai sostenitori di quel comunismo elementare, ovvero l’espropriazione della miseria.

Alle grevi teorie comuniste, Marx contrappone l’ideale di un comunismo inteso come superamento attivo della proprietà privata. Tale comunismo, in quanto processo storico, salvaguarda l’eredità di tutto il pregresso sviluppo, preservandone le conquiste sotto forma di competenze, necessità, patrimonio materiale e immateriale dell’umanità. Il comunismo rende questo patrimonio accessibile a ciascun membro della società. In quanto divenire storico, esso si caratterizza per la sua natura cosciente e non spontanea, poiché la socializzazione implica una relazione consapevole dei membri della società con i processi di produzione e riproduzione della vita sociale. Il comunismo sottende l’appropriazione delle leggi oggettive che governano l’evoluzione sociale, e l’umanità può metterle al proprio servizio, così come già fa con le leggi della fisica, della chimica e della biologia nell’ambito dell’industria moderna.

L’intera storia umana precedente prepara il terreno per un tipo di comunismo che trascende la proprietà privata, pur preservandone l’eredità e i frutti del progresso compiuto. Parallelamente, si sviluppa una coscienza adeguata, ovvero la teoria del comunismo. Il comunismo si configura come il processo di superamento della proprietà privata: «Per abbattere il concetto di proprietà privata, è sufficiente l’idea stessa del comunismo. Per sradicare la proprietà privata nella realtà, è necessaria un’azione comunista concreta e tangibile».

È doveroso evidenziare un aspetto che, a nostro avviso, riveste particolare importanza. Ci riferiamo all'impiego delle espressioni "socialismo" e "comunismo", il cui significato, come comunemente inteso, si è definito solo progressivamente. Negli scritti del 1844, a cui abbiamo fatto precedentemente riferimento, il socialismo viene ancora descritto come una società nella quale la proprietà privata non è semplicemente abolita, ma è già stata soppressa. Il comunismo, invece, è concepito come il processo – o meglio, l'atto – di negazione (cioè di eliminazione) della proprietà privata. Per tale ragione, in questo specifico contesto, esso non è considerato una forma di organizzazione sociale, bensì unicamente come un principio ispiratore, una forza propulsiva rivolta al futuro immediato.

Questa particolarità permette di comprendere appieno le ragioni per cui Marx ed Engels, appena tre anni dopo la stesura del “Manifesto del Partito Comunista”, affrontarono e presentarono tali questioni nelle “Manoscritte”. In esse, si possono discernere con maggiore chiarezza le origini del principio cardine del “Manifesto”: «La proprietà privata borghese, nella sua forma moderna, rappresenta l’ultima e più compiuta espressione di un sistema di produzione e di appropriazione dei beni fondato sull’antagonismo di classe e sullo sfruttamento reciproco. Pertanto, i comunisti possono riassumere la loro teoria in questa singolare affermazione: abolizione della proprietà privata».

Tale abolizione, come già anticipato, implica la salvaguardia dell’intero patrimonio culturale, delle capacità produttive e dei bisogni che si sono sviluppati sulla base della proprietà privata. Essa non è fine a sé stessa, ma è strumentale alla realizzazione di una società comunista dotata di un sistema di proprietà e di comunicazione adeguati.

Qualunque forma assuma la società comunista, l’uomo, in quanto essere naturale, dovrà vivere in armonia con la natura e appropriarsene. Ogni elemento di cui l’uomo si serve e che consuma è costituito da materiale naturale, trasformato attraverso processi naturali posti al servizio dell’uomo, oppure è direttamente costituito da tali processi stessi. Vent’anni dopo la pubblicazione del «Manifesto», nell’Introduzione degli «Scritti economici del 1857–1859», Marx scriveva: «Ogni forma di produzione implica l’appropriazione, da parte dell’individuo, di elementi provenienti dalla natura, entro i limiti di una determinata struttura sociale e mediante essa. In tal senso, sarebbe una mera tautologia affermare che la proprietà (l’appropriazione) è una condizione della produzione. Sarebbe però assurdo concludere da ciò che si possa stabilire una specifica forma di proprietà, come ad esempio la proprietà privata – la quale, peraltro, presuppone come condizione preliminare la sua antitesi: l’assenza di proprietà».

In una società comunista, la proprietà deve assumere una forma collettiva, ovvero pubblica, per quanto concerne tutti i mezzi di produzione sociale. Di conseguenza, la gestione dell’intera produzione sociale dovrà essere anch’essa pubblica e unitaria. Allo sviluppo del carattere pubblico della produzione, fondato sulla proprietà privata, dovrà corrispondere un analogo carattere pubblico nell’appropriazione dei suoi frutti. Marx ed Engels ribadirono questo principio in numerose opere. Forse, questa è l’unica tesi che i comunisti contemporanei hanno appreso a memoria, sebbene non ne afferrino appieno il significato.

Al fine di dare avvio a questo processo storico nei paesi più avanzati, Marx ed Engels proposero le seguenti misure: «1. Espropriazione della proprietà terriera e destinazione dei relativi proventi alla copertura delle spese statali.»

2. L’introduzione di un’imposta progressiva particolarmente elevata. 3. L’abolizione dell’imposta di successione. 4. La confisca dei beni di tutti gli emigranti e dei ribelli. 5. La centralizzazione del sistema creditizio nelle mani dello Stato, mediante la creazione di una banca nazionale dotata di capitale pubblico e investita di un monopolio esclusivo. 6. La concentrazione di tutti i servizi di trasporto sotto il controllo statale. 7. L’incremento del numero di fabbriche statali, l’implementazione di strumenti di produzione, la bonifica dei terreni e il miglioramento delle terre, il tutto secondo un piano generale. 8. L’obbligo di lavoro per tutti i cittadini, attraverso la creazione di corpi di lavoratori organizzati, in particolare nel settore agricolo. 9. L’integrazione tra agricoltura e industria e la promozione di una graduale riduzione delle disparità tra aree urbane e rurali. 10. Un sistema di istruzione pubblica e gratuita per tutti i bambini, con l’eliminazione del lavoro minorile nelle fabbriche nella sua forma attuale.

Come si può evincere, alla luce dell’obiettivo di superare la proprietà privata, queste misure non appaiono eccessivamente radicali, ma nel loro insieme si dimostrano coerenti, come del resto ci si aspetterebbe da un manifesto. Un manifesto deve, infatti, essere chiaro, comprensibile e delineare la direzione principale dell’evoluzione. Inoltre, gli stessi autori riconoscevano che tali misure sono subordinate alle condizioni storico-concrete e, di conseguenza, possono variare. Ci soffermeremo specificamente su di esse in uno dei prossimi articoli, analizzandole alla luce dell’esperienza del socialismo nel corso del XX secolo.

Aggiungiamo ora alcune considerazioni circa la natura della proprietà privata, e su cosa, in particolare, i comunisti intendano eliminare, preservando al contempo tutti gli elementi positivi, affinché «al posto della vecchia società borghese, con le sue classi e le opposizioni di classe», si affermi «un’associazione in cui il libero sviluppo di ciascuno sia la condizione per il libero sviluppo di tutti».

La proprietà privata scaturisce da un'attività lavorativa priva di autonomia e, al contempo, costituisce lo strumento per generare un tipo di lavoro altrettanto alienato: un'attività in cui né il processo né il risultato finale appartengono direttamente alla comunità dei lavoratori, un'attività che, intrinsecamente, non coinvolge i lavoratori, ma che si configura unicamente come mezzo per il conseguimento di fini esterni (il profitto). L'essenza della proprietà privata risiede in questa peculiare forma di lavoro. Tale forma di lavoro raggiunge il suo apice di sviluppo nel sistema capitalista, attraverso una produzione di merci su larga scala, resa possibile da un'estrema parcellizzazione del lavoro, che trasforma i frutti dell'attività umana in merce

Per lungo tempo, la divisione sociale del lavoro ha contribuito al progresso dell’umanità, incrementando la produttività e ponendo le fondamenta materiali per l’automatizzazione di compiti ripetitivi, attività che l’uomo, a causa della limitata evoluzione degli strumenti a sua disposizione, era costretto a svolgere. Ma non si tratta semplicemente di suddividere le mansioni, assegnando a ciascuno un ruolo specifico – Vasja meccanico, Petya programmatore – e di affidare la supervisione del loro lavoro a un responsabile, il tutto in nome di un presunto bene comune. In realtà, nessuno dei partecipanti a questo processo persegue attivamente tale obiettivo. In questa dinamica, gli individui sono legati tra loro in modo indiretto, attraverso i prodotti del loro lavoro. Sono obbligati a scambiarsi i risultati della propria attività con quelli di cui necessitano, mantenendo un atteggiamento di indifferenza reciproca durante il processo produttivo. Non importa chi e come realizzi un determinato prodotto; ciò che conta è che quel prodotto, frutto del lavoro altrui, sia utile a me.

L’evoluzione della divisione del lavoro all’interno della società ha attraversato diverse tappe fondamentali, che hanno avuto un impatto positivo sull’avanzamento della civiltà umana. Sarebbe auspicabile dedicare un’opera a parte all’analisi approfondita di questo processo. In questo contesto, è sufficiente sottolineare che il progresso sociale è il frutto della storia. I marxisti rifiutano con fermezza ogni ipotesi riguardante una presunta predisposizione biologica o innata dei membri della società verso determinate attività lavorative. Tali affermazioni non superano il vaglio di una critica scientifica anche solo superficiale, poiché non solo contrastano con quanto emerso dagli studi storici sull’argomento, ma anche con le dinamiche che osserviamo quotidianamente. Ogni anno, infatti, nuove professioni emergono, mentre altre scompaiono. I genitori spesso faticano a comprendere la natura dell’attività svolta dai propri figli e, di conseguenza, non potrebbero trasmettere loro per via ereditaria, insieme ai geni, un’attitudine specifica per tale lavoro.

Da quando il lavoro si è frammentato e i frutti dell’attività umana sono stati ridotti a merci, ovvero concepiti primariamente come oggetti di scambio, questi prodotti hanno iniziato a esercitare una forma di dominio sull’uomo. Le cose, i manufatti nati dall’ingegno umano, e il denaro, in quanto simbolo di ogni bene materiale, hanno acquisito un potere innegabile sulle persone. E, poiché la forza lavoro è divenuta essa stessa merce, questo potere si è concentrato nelle mani di coloro che dispongono del capitale necessario per acquistare i mezzi di produzione e la manodopera. Queste figure, in sostanza, detengono il controllo sul lavoro altrui. Dal punto di vista strettamente economico, essi rappresentano semplicemente il capitale. La loro funzione primaria è quella di accrescere tale capitale, sfruttando il lavoro di altri, in modo che il valore complessivo aumenti progressivamente. In ultima analisi, questo meccanismo conduce alla condizione in cui la maggior parte dei membri della società, in un modo o nell’altro, contribuisce all’espansione del capitale, alimentandone la crescita esponenziale.«Unicamente i prodotti del lavoro privato, svolto in autonomia e indipendenza reciproca, si presentano come merci in opposizione l’uno all’altro». Di conseguenza, l’abolizione della proprietà privata implica l’abolizione della produzione di merci in senso generale, e non soltanto della produzione capitalistica di merci.

Il valore tende intrinsecamente all’auto-espansione. Questa è la legge che regola una società in cui ogni individuo è legato alla propria attività parziale; in altre parole, una società caratterizzata da una produzione di merci sviluppata. Le leggi degli stati moderni costituiscono l’espressione giuridica di tale legge oggettiva della produzione, fondata su un lavoro non libero e parcellizzato.

Fino a un certo punto, l’impulso intrinseco del capitale alla sua stessa crescita rappresenta la forza motrice dello sviluppo sociale. Tuttavia, esso genera forze sociali che si sentono soffocate entro i limiti imposti dal capitale. Il loro impiego esclusivo al fine di accrescere il capitale conduce a conseguenze catastrofiche.

È imperativo governare tali dinamiche, e, anzi, è possibile farlo a vantaggio dell’intera società unicamente attraverso una collaborazione sinergica, volta al bene comune, guidata da un piano organico e razionale. A tal fine, è imprescindibile coltivare le capacità umane universali – affinare la sensibilità, promuovere un pensiero e una moralità inclusivi – trascendendo i confini ristretti della propria professione, che ostacolano persino la comprensione delle discipline affini, figuriamoci la percezione dei processi sociali nella loro complessità. Di conseguenza, l’abolizione della proprietà privata equivale all’eliminazione di un’attività lavorativa priva di autonomia. Questo processo si compie progressivamente con lo sviluppo dell’attività comunista, liberandola e trasformandola in un bene a disposizione di tutti gli esseri umani. Il comunismo si contrappone alla proprietà privata, proponendo la proprietà comune come forma di relazione pratica e collettiva tra gli individui, i mezzi di produzione e le risorse naturali.

La proprietà collettiva assicura un impiego più efficace delle risorse produttive, a beneficio non solo dell’intera società, ma anche di ciascuno dei suoi membri. Essa presuppone un coinvolgimento attivo e una partecipazione consapevole di ogni singolo individuo rispetto al patrimonio comune. Pertanto, se la proprietà privata è il risultato della divisione del lavoro – e della sua progressiva frammentazione in mansioni sempre più specializzate \footnote{Proprio perché nell'URSS la divisione sociale del lavoro non fu superata (e la tendenza a superarla non divenne prevalente), e il proletariato continuò a esistere in quanto tale, ovvero come classe appartenente alla società borghese, si assistette a una completa restaurazione del dominio della proprietà privata} – allora la proprietà collettiva implica, come già evidenziato, la socializzazione del lavoro e la superamento della sua parcellizzazione. La proprietà pubblica promuove un’attività umana libera, fondata sui principi di verità, bene e bellezza. Quest’ultimo aspetto costituisce il vero fondamento economico per una collettivizzazione autentica, e non meramente apparente, e per una reale convergenza tra gli interessi individuali e quelli collettivi di tutti i membri della società

La proprietà collettiva o sociale non andrebbe concepita come la proprietà di aziende agricole o cooperative, o di altri consorzi di individui che producono beni per il mercato. Queste entità, infatti, non trascendono i limiti della proprietà privata, poiché rimangono confinate nell'ambito della produzione di merci. I membri di tali consorzi sono legati alla società in quanto detentori di merci e produttori, e la natura stessa della produzione di merci definisce le dinamiche interne, ponendo al centro la creazione di valore e facendo del profitto lo scopo ultimo dell'attività collettiva.

Non varrebbe la pena cercare le origini della proprietà collettiva, per esempio, nella proprietà familiare. Nel XIX secolo, alcuni socialisti ipotizzarono che la proprietà sociale potesse nascere da una sorta di “estensione” della proprietà familiare a una scala più ampia, sociale. La società veniva concepita come una grande famiglia, con i metodi di distribuzione tipici di un nucleo familiare. Ma la famiglia è l’incarnazione stessa del carattere privato dell’appropriazione dei prodotti del lavoro. In quanto tale, nella sua forma moderna, essa è il risultato dello sviluppo della proprietà privata. Laddove la famiglia organizza la produzione e la riproduzione dell’uomo in quanto forza lavoro – l’individuo come agente economico all’interno di una società di produzione di merci – essa è interamente integrata nel sistema e si riproduce pienamente secondo le leggi che governano il movimento della proprietà privata.Poiché essa riproduce l’intera sua compagine e costituisce l’unità di consumo primaria, la sua essenza è interamente plasmata dalla logica della produzione capitalistica. Tale implicazione si manifesta inequivocabilmente nelle dinamiche salariali e nella generazione del profitto. Inoltre, la famiglia moderna, in quanto entità economica autonoma, non solo non esclude, ma incorpora attivamente lo sfruttamento reciproco tra gli individui, un principio intrinsecamente incompatibile con l’ideale comunista. Di conseguenza, la proprietà familiare non può essere legittimamente considerata il fondamento di una proprietà sociale, che rappresenterebbe l’architrave di un sistema comunista.

E men che meno la proprietà comunista, intesa come punto di partenza per la transizione verso una società comunista, può essere identificata in qualsiasi forma di comunità isolata, estranea alla rete della produzione sociale (ad esempio, comunità religiose basate su un’economia di sussistenza).

Ma allora, in quali contesti è possibile osservare la proprietà comune, elemento distintivo del comunismo? È evidente che ciò si verifica unicamente in quelle situazioni in cui l’attività umana si svolge in forma direttamente collettiva e libera. Tale attività, analizzata sotto il profilo dell’impiego e della gestione delle risorse, nonché dell’appropriazione attraverso forme sociali, costituisce la manifestazione concreta della proprietà comune.

Nel capitolo precedente, abbiamo fornito solo alcuni esempi di questa attività comunista, sia nella società socialista che in quella capitalista. Se esaminiamo tali esempi, e altri, alla luce dell’espressione della volontà in relazione all’oggetto del lavoro, ai mezzi di produzione e al prodotto finale, potremo individuare le diverse forme della proprietà sociale e il carattere intrinsecamente sociale dell’appropriazione.

Un database pubblico, per sua stessa essenza, non può più essere considerato una proprietà privata, men che meno un elemento di un sistema economico centralizzato. Lo stesso principio si applica al software libero e ai molteplici servizi che si fondano sulla produzione e sull’utilizzo condivisi. Esistono, inoltre, server che mettono in comune le proprie risorse, un modello il cui sfruttamento basato sulla proprietà privata risulterebbe inefficiente e controproducente, come nel caso delle biblioteche e degli archivi digitali dedicati a film e musica, dove gli utenti contribuiscono attivamente, offrendo e ricevendo contenuti. Tuttavia, l’esempio del software e dei database è più calzante, poiché sottolinea maggiormente l’aspetto della produzione rispetto a quello del mero consumo.

Da quanto precede, si desume che la proprietà pubblica scaturisce dall'effettivo dinamismo produttivo, anche in quelle situazioni in cui non si è ancora giunti a una sua formalizzazione nell'ambito della socializzazione dei mezzi di produzione, all'interno del sistema capitalista. La formalizzazione di tale socializzazione – ovvero la conquista del potere politico da parte del proletariato e la conseguente appropriazione dei mezzi di produzione – aprirà la strada a un più ampio sviluppo della reale proprietà pubblica. È importante sottolineare che, laddove esista una proprietà pubblica effettiva, e non meramente formale, verrà meno ogni pretesto per tentare di appropriarsi indebitamente di una parte delle risorse pubbliche, oppure, al contrario, per disperderle, gestendole con superficialità come se non appartenessero a nessuno.

Qualora si optasse per una collettivizzazione di natura formale, risulterebbe imprescindibile una regolamentazione esterna, se non altro sotto forma di un'elevata coscienza morale (in realtà, tradotta in leggi e in un sistema di coercizione atto a garantirne l'applicazione) concernente gli strumenti, l'oggetto e i frutti del lavoro, escludendo questi ultimi dalla collettivizzazione. In caso di una collettivizzazione autentica, la necessità di tale regolamentazione decadrebbe, poiché la cultura riguardante i mezzi di produzione diverrebbe semplicemente una premessa per il loro impiego
\end{document}
